\section{Objetivo general}
Diseñar e implementar una librería npm que integre un pipeline de ingeniería de características y un modelo híbrido de aprendizaje automático capaz de detectar amenazas conocidas y anomalías estadísticas en endpoints Node.js en tiempo de ejecución, para dar a startups sin equipo de seguridad una herramienta accesible de protección a nivel de capa de aplicación sin comprometer el rendimiento.

\section{Objetivos específicos}
\begin{itemize}
\item Implementar un pipeline de ingeniería de características que se ejecute de forma asíncrona mediante \textit{worker\_threads}, con una latencia adicional máxima de 5 ms por petición.
\item Desarrollar un modelo híbrido compuesto por un Random Forest y un Isolation Forest, alcanzando un F1-score mínimo de 0.80 en al menos 3 de las 4 categorías de ataque para el Random Forest, y un recall mínimo del 50\% con un máximo de 10\% de falsos positivos para el Isolation Forest, ambos evaluados sobre un conjunto de prueba bloqueado.
\item Lograr una detección mínima del 80\% de los payloads de ataque por categoría en un entorno controlado, manteniendo un sobrecosto de latencia media del servidor que no exceda los 15 ms respecto a la línea base, o del 10\% en entornos de producción representativos con latencia base mayor o igual a 100 ms.
\end{itemize}
